\section{Status, scorecard and open problems}
\label{sec:status}

\begin{table}[H]
  \centering
  \small
  \begin{tabular}{L{1.2cm}L{4.4cm}L{6.6cm}L{1.6cm}}
    \toprule
    & question & answer so far & status \\
    \midrule
    M0 & Are the noisy numbers real? & Not before M0. Floor now \VPostFinalMin{} to
      \VPostFinalMax{}; a bias, not a convergence failure (\cref{sec:noise}). & done \\
    M1 & Which gradient gives most precision per shot? & RVB: ADAM + parameter shift
      (\VMIdealRvbPs). HVA: plain SPSA (\VMIdealHvaSpsa). & done \\
    M2 & How much bias can mitigation remove? & -- & next \\
    M3 & Is the benchmark testing anything? & Not much: exact at zero parameters
      (\cref{thm:cycle}). & planned \\
    M4 & How far can depth fall? & Free first step found: bind parameters before
      transpiling, \VCxHva{} $\to$ \VBoundCxHva{} CX (\cref{prop:bound}). & planned \\
    M5 & Can subspace methods remove the rest? & -- & planned \\
    \bottomrule
  \end{tabular}
  \caption{Running scorecard, one row per roadmap milestone (\py{docs/roadmap.md}).}
  \label{tab:scorecard}
\end{table}

\subsection{Open problems}

\begin{enumerate}
  \item \textbf{The noisy floor} (\VFloorHvaFull{} at the exact parameters) is the dominant
    error, about \VPostErrHvaSpsa. Better optimisers cannot remove it (\cref{prop:mixed} and \cref{prop:depol}).
    The near-multiplicative structure of Table~\ref{tab:floor} favours rescaling and
    extrapolation \parencite{temme2017} (M2) and fewer CX (M4).
  \item \textbf{Hardware calibration H0} is deferred; every noisy number here is a simulation of
    a 2022 calibration snapshot.
  \item \textbf{M1 follow-ups:} decaying ADAM learning rate for HVA; longer RVB runs; more seeds
    with device noise.
  \item \textbf{The instance is easy} (\cref{rem:easy}); results should be re-validated on M3
    instances before conclusions generalise.
\end{enumerate}

\subsection{Issues found while writing this report}

These were found by recomputing every number from the code. The development session
reproduced each of them and is fixing the prose and code on \py{main}.
\begin{itemize}
  \item Parametric transpilation leaves each Heisenberg block at 5 CX instead of 3
    (\cref{prop:bound}). Binding first gives \VBoundCxHva{} CX and keeps \VBoundKeepHva{} instead of
    \VParamKeepHva{} of the energy at a generic point.
  \item \py{docs/physics.md} quotes 93 CX for RVB; the code gives \VCxRvb.
  \item The README's seed spreads were population standard deviations (Table~\ref{tab:m1},
    footnote); \py{main} is switching aggregated tables to sample standard deviations.
  \item The M0 spread figures and those in this report were measured under different conditions,
    and both are valid. The roadmap's figures (\VRoadmapSpreadLow{} at 1024 shots, \VRoadmapSpreadHigh{} at 4096) were taken
    \emph{with} \py{FakeGuadalupeV2} noise at the exact parameters, where noise breaks the
    cancellation of \cref{prop:zerovar}. This report's figures (\VSpreadLowReal, \VSpreadHighReal)
    are noise-free at the start point. Without noise the spread at the optimum is exactly zero,
    so the roadmap now states the conditions of its measurement.
  \item \py{HVAnsatz}'s docstring writes $e^{-i\gamma\Hev}e^{-i\phi\Hod}$; the circuit implements
    $e^{-i\frac\gamma2\Hev}e^{-i\frac\phi2\Hod}$ (\cref{dfn:ansatz}). Only the convention differs.
\end{itemize}
